%==============================================================================
%  RESUME — TRACK B  (research/projects-led)
%
%  Section order: Education -> Technical Skills -> Projects
%                 -> Professional Experience
%  Use for: ML/AI Engineering and Embedded/Test Automation, where the work
%  itself is more persuasive than the employer's name.
%
%  Identical preamble to Track A on purpose — only the body order differs.
%
%  RULES FOR THE TAILORING AGENT
%   - ONE PAGE. Not a preference. If it spills, cut bullets, never shrink margins.
%   - Never change the document class, margins, or the \newcommand definitions.
%   - Only content between \begin{document} and \end{document} is rewritten.
%   - Delete unused blocks entirely; never leave a {{TOKEN}} in a sent resume.
%   - Escape % & _ # $ in content ( 94\% , R\&D ).
%   - Compiles with Tectonic locally AND on Overleaf unchanged.
%==============================================================================

\documentclass[letterpaper,10pt]{article}

\usepackage{latexsym}
\usepackage[empty]{fullpage}
\usepackage{titlesec}
\usepackage{marvosym}
\usepackage[usenames,dvipsnames]{color}
\usepackage{verbatim}
\usepackage{enumitem}
\usepackage{xcolor}
\usepackage[
    colorlinks=true,
    linkcolor=black,
    urlcolor=blue,
    citecolor=blue
]{hyperref}
\usepackage{fancyhdr}
\usepackage[english]{babel}
\usepackage{tabularx}
% ATS text extraction: glyphtounicode is a pdfTeX feature. Guard it so the file
% compiles unchanged under pdfLaTeX (Overleaf's default) AND XeTeX (Tectonic),
% where correct ToUnicode mapping is produced natively anyway.
\ifdefined\pdfglyphtounicode
  \IfFileExists{glyphtounicode.tex}{\input{glyphtounicode}}{}
\fi
\pagestyle{fancy}
\fancyhf{}
\fancyfoot{}
\renewcommand{\headrulewidth}{0pt}
\renewcommand{\footrulewidth}{0pt}

\addtolength{\oddsidemargin}{-0.6in}
\addtolength{\evensidemargin}{-0.6in}
\addtolength{\textwidth}{1.2in}
\addtolength{\topmargin}{-0.85in}
\addtolength{\textheight}{1.5in}

\urlstyle{same}
\raggedbottom
\raggedright
\setlength{\tabcolsep}{0in}
\setlength{\parskip}{0pt}

\setlist[itemize]{topsep=0pt, itemsep=0.4pt, parsep=0pt, leftmargin=0.15in}

\titleformat{\section}{
  \vspace{2pt}\scshape\raggedright\normalsize\bfseries
}{}{0em}{}[\color{black}\titlerule \vspace{2pt}]

\ifdefined\pdfgentounicode \pdfgentounicode=1 \fi

%-------------------------
% Custom commands
%-------------------------

\newcommand{\resumeSubheading}[4]{
  \vspace{1pt}\item
    \begin{tabular*}{0.985\textwidth}[t]{l@{\extracolsep{\fill}}r}
      \textbf{#1} & #2 \\
      \textit{\small#3} & \textit{\small #4} \\
    \end{tabular*}\vspace{1pt}
}

\newcommand{\resumeSingleHeading}[2]{
  \vspace{1pt}\item
    \begin{tabular*}{0.985\textwidth}{l@{\extracolsep{\fill}}r}
      \textbf{#1} & \small#2 \\
    \end{tabular*}\vspace{2pt}
}

% Sub-project line inside a research role, e.g. "Project: MiGa - Multi-Chicken Gait Assessment"
\newcommand{\resumeProjectLine}[1]{
  \vspace{1pt}\item[]{\small\textit{#1}}\vspace{2pt}
}

\newcommand{\resumeItem}[1]{\item\small{#1}}

\newcommand{\resumeSubHeadingListStart}{\begin{itemize}[leftmargin=0.0in, label={}]}
\newcommand{\resumeSubHeadingListEnd}{\end{itemize}\vspace{-2pt}}
\newcommand{\resumeItemListStart}{\begin{itemize}[leftmargin=0.17in]}
\newcommand{\resumeItemListEnd}{\end{itemize}\vspace{2pt}}

%==============================================================================
\begin{document}

%----------HEADING----------
% Full URLs, not hyperlinked words: ATS text extraction keeps them.
\begin{center}
    \textbf{\Large \scshape {{FULL_NAME}}} \\ \vspace{2pt}
    \small {{PHONE}} $|$
    \href{mailto:{{EMAIL}}}{\textcolor{blue}{\underline{{{EMAIL}}}}} $|$
    \href{{{PORTFOLIO_URL}}}{\textcolor{blue}{\underline{{{PORTFOLIO_URL}}}}} $|$
    \href{{{GITHUB_URL}}}{\textcolor{blue}{\underline{{{GITHUB_URL}}}}}
\end{center}
\vspace{-6pt}

%-----------EDUCATION-----------
% Track B keeps the bachelor's more often than Track A — for Embedded roles the
% Electronics & Communication degree is directly relevant, not filler.
\section{Education}
  \resumeSubHeadingListStart
    \resumeSubheading
      {{{INSTITUTION_1}}}{{{EDU_DATES_1}}}
      {{{DEGREE_1}}}{}
    \resumeSubheading
      {{{INSTITUTION_2}}}{{{EDU_DATES_2}}}
      {{{DEGREE_2}}}{}
  \resumeSubHeadingListEnd

%-----------TECHNICAL SKILLS-----------
% Category presets per track are in system/profile/skills-matrix.md.
% Reorder so the JD's top requirements land first. 4-6 lines.
\section{Technical Skills}
\begin{itemize}[leftmargin=0.0in, label={}]
    \small{\item{
        \textbf{ {{SKILL_CATEGORY_1}} }{: {{SKILL_LIST_1}} } \\
        \textbf{ {{SKILL_CATEGORY_2}} }{: {{SKILL_LIST_2}} } \\
        \textbf{ {{SKILL_CATEGORY_3}} }{: {{SKILL_LIST_3}} } \\
        \textbf{ {{SKILL_CATEGORY_4}} }{: {{SKILL_LIST_4}} } \\
        \textbf{ {{SKILL_CATEGORY_5}} }{: {{SKILL_LIST_5}} }
    }}
\end{itemize}
\vspace{-4pt}

%-----------PROJECTS-----------
% Leads the page on this track. The FIRST project is the signature project —
% chosen for this employer alone, never reused. See system/data/signature-projects.md.
\section{Projects}
  \resumeSubHeadingListStart

    \resumeSingleHeading{{{PROJECT_1_NAME}}}{{{PROJECT_1_STACK}}}
      \resumeItemListStart
        \resumeItem{{{PROJECT_1_BULLET_1}}}
        \resumeItem{{{PROJECT_1_BULLET_2}}}
        \resumeItem{{{PROJECT_1_BULLET_3}}}
      \resumeItemListEnd

    \resumeSingleHeading{{{PROJECT_2_NAME}}}{{{PROJECT_2_STACK}}}
      \resumeItemListStart
        \resumeItem{{{PROJECT_2_BULLET_1}}}
        \resumeItem{{{PROJECT_2_BULLET_2}}}
      \resumeItemListEnd

  \resumeSubHeadingListEnd

%-----------PROFESSIONAL EXPERIENCE-----------
% Research roles use \resumeProjectLine for their named sub-projects, the way
% her _A / _DA / _SA / _SDA variants do.
\section{Professional Experience}
  \resumeSubHeadingListStart

    \resumeSingleHeading
      {{{ROLE_1}}, {{EMPLOYER_1}}}{{{EMP_DATES_1}}}
      \resumeProjectLine{Project: {{EMP_1_SUBPROJECT_1}}}
      \resumeItemListStart
        \resumeItem{{{EMP_BULLET_1_1}}}
        \resumeItem{{{EMP_BULLET_1_2}}}
      \resumeItemListEnd
      \resumeProjectLine{Project: {{EMP_1_SUBPROJECT_2}}}
      \resumeItemListStart
        \resumeItem{{{EMP_BULLET_1_3}}}
      \resumeItemListEnd

    \resumeSingleHeading
      {{{ROLE_2}}, {{EMPLOYER_2}}}{{{EMP_DATES_2}}}
      \resumeItemListStart
        \resumeItem{{{EMP_BULLET_2_1}}}
        \resumeItem{{{EMP_BULLET_2_2}}}
        \resumeItem{{{EMP_BULLET_2_3}}}
      \resumeItemListEnd

  \resumeSubHeadingListEnd

\end{document}
